% Figure: even/odd symmetry of integrands. An odd integrand over a
% symmetric interval integrates to zero (the two signed areas cancel);
% an even integrand integrates to twice the half-interval. Left panel:
% odd integrand x^3 on [-1.4, 1.4]; right panel: even integrand x^2 on
% the same interval.
\begin{figure}[htbp]
\centering
\begin{tikzpicture}

\begin{axis}[
  name=odd,
  width=6.0cm, height=5.2cm,
  xlabel={$x$}, ylabel={$f(x)$},
  xmin=-1.5, xmax=1.5, ymin=-3, ymax=3,
  domain=-1.4:1.4, samples=120,
  axis lines=middle, font=\footnotesize,
  title={Odd: $f(x)=x^{3}$},
]
\addplot[draw=none, fill=engred!15, domain=-1.4:0] {x^3} \closedcycle;
\addplot[draw=none, fill=engblue!15, domain=0:1.4] {x^3} \closedcycle;
\addplot[engblack, very thick] {x^3};
\node[font=\scriptsize, engred] at (axis cs:-0.85,-1.5) {$-$};
\node[font=\scriptsize, engblue] at (axis cs:0.85,1.5) {$+$};
\end{axis}

\begin{axis}[
  at={(odd.east)}, anchor=west, xshift=16pt,
  name=even,
  width=6.0cm, height=5.2cm,
  xlabel={$x$}, ylabel={},
  xmin=-1.5, xmax=1.5, ymin=-0.3, ymax=2.2,
  domain=-1.4:1.4, samples=120,
  axis lines=middle, font=\footnotesize,
  title={Even: $f(x)=x^{2}$},
]
\addplot[draw=none, fill=engblue!15, domain=-1.4:1.4] {x^2} \closedcycle;
\addplot[engblack, very thick] {x^2};
\node[font=\scriptsize, engblue] at (axis cs:-0.75,0.9) {$+$};
\node[font=\scriptsize, engblue] at (axis cs:0.75,0.9) {$+$};
\end{axis}

\end{tikzpicture}
\caption[Even and odd symmetry of an integrand]{Symmetry read off the
integrand before any antiderivative is found. An odd integrand (left,
$x^{3}$) over an interval symmetric about the origin contributes a
negative signed area on the left half and an equal positive signed area
on the right half, so the two cancel and the integral is zero. An even
integrand (right, $x^{2}$) contributes equal positive areas on each
half, so the integral is twice the half-interval value. Reading the
symmetry first turns $\int_{-a}^{a} x^{3}\dd x = 0$ and
$\int_{-a}^{a} x^{2}\dd x = 2\int_{0}^{a} x^{2}\dd x$ into one-line
results and guards against a sign slip in the antiderivative.}
\label{fig:vol02ch06:even-odd-symmetry}
\end{figure}
